% s-113-p113-m003; original PDF 113, 114
\Aphorism{s-113-p113-m003}{128}
Ich sehe einen hauptsächlichen Rangunterschied der Weiber in der Art, wie sie ihren Körper zu verrathen wissen. Hier ist, zwischen den Grenzen. Exhibition und „Faltenhemd“, Aphrodite und Artemis, fast alles möglich: Blick, Stimmklang, Gebärde, Gang, Kleidung — alle diese Sprachen reden von dem Einen, alle diese Schleier sind ein verhüllendes Sehenlassen des Einen. Übrigens ein Fall, wo die Übersetzung mehr verspricht, als das Original halten kann; die Süssigkeit liegt in der Lockung, nicht in dem, wozu gelockt wird. Die ersten vier Acte spannen, sind geniessbar; der fünfte Act, der eigentliche Act, fällt ab. — Wie lebt eigentlich ein junges Weib? Ein Mann kann sich eher in die Existenz einer Eidechse hineindenken als in die seltsame Unwahrscheinlichkeit des weiblichen Daseins, diese schamhaft-schamlose Mischung von Baubo und Urania. So als Räthsel mit allbekannter Lösung herumlaufen müssen, ewig nur \Foreign{sub specie} desselben „höheren Zwecks“ in Betracht kommen, nur als prächtige Instrumentation einer sehr kurzen, sehr trivialen Melodie gehört werden! Mich wundert, dass dem Weibe dabei nicht eine tiefe Verachtung des Mannes und seiner selbst anfliegt. Und dieses Eine, als dessen Drum und Dran das junge Weib sich selber fühlt, äusserlich und psychisch so erschwert: ein ganzer Berg von Moral und Romantik um diese lächerliche Maus!
\par

